\section{Comprensión intuitiva y mejoras}

\subsection{Aprendizaje por imitación ponderado por reward}

El gradiente de política es en realidad \textbf{aprendizaje por imitación ponderado por el reward de la trayectoria}: refuerza las acciones con reward alto y suprime las acciones con reward bajo. A diferencia de la supervisión pura, esta formulación nos permite explorar mediante una stochastic policy.

\begin{knowledgebox}{La naturaleza del aprendizaje por ensayo y error}
\begin{itemize}
    \item la política genera las muestras, y el reward determina la dirección del gradient;
    \item maneja el reward disperso con más flexibilidad, porque el muestreo incluye la trayectoria completa;
    \item el objetivo de optimización es el reward esperado, no la exactitud de un solo paso.
\end{itemize}
\end{knowledgebox}

\subsection{Reward-to-Go y causalidad}

Queremos que el gradient de un time step determinado solo se vea afectado por el reward de ese step y los posteriores. Usar el reward-to-go permite lograr la causalidad:
$$\nabla_\theta J(\theta) \approx \frac{1}{N}\sum_{i=1}^{N}\sum_{t=1}^{T} \nabla_\theta \log \pi_\theta(a_{i,t} \mid s_{i,t}) \left(\sum_{t'=t}^{T} r(s_{i,t'}, a_{i,t'}) \right).$$
Esta expresión, en cada acción, solo observa el reward futuro, por lo que se ajusta mejor al causal effect.

\begin{table}[H]
\centering
\small
\begin{tabular}{p{0.3\textwidth}p{0.45\textwidth}p{0.2\textwidth}}
\toprule
Término & Significado & Función \\
\midrule
Trajectory reward & reward acumulado de toda la trayectoria & supervision global \\
Reward-to-go & reward desde el instante actual hasta el final & mantiene la causalidad \\
Baseline & una constante o una función de valor & reduce la variance \\
\bottomrule
\end{tabular}
\caption{Expresiones de reward comunes en la estimación del gradiente de política}
\end{table}

\subsection{Baseline y Advantage}

\begin{warningbox}{Sensible a la escala del reward}
Si todos los rewards son positivos, el gradiente aumenta la verosimilitud de todas las acciones, incluso si algunas apenas son «no tan malas». Esto incrementa considerablemente la variance del gradiente.
\end{warningbox}

Al restar un baseline $b(s_t)$ del reward-to-go, no cambiamos la esperanza, pero podemos reducir la variance:
$$\nabla_\theta J(\theta) = \mathbb{E}_{\tau \sim p_\theta(\tau)} \left[\sum_t \nabla_\theta \log \pi_\theta(a_t \mid s_t) \cdot (G_t - b(s_t))\right],$$
donde $G_t$ es el reward-to-go. Si $b(s_t)=V^\pi(s_t)$, se obtiene el advantage estimator $A^\pi(s_t,a_t)$.

\begin{importantbox}{La intuición del Advantage}
El Advantage mide en qué medida una acción es mejor que el promedio del estado actual. Un advantage positivo refuerza esa acción, uno negativo la suprime.
\end{importantbox}

\subsection{Objetivo sustituto y diferenciación automática}

Al incorporar el reward-to-go y el baseline en el surrogate objective:
$$\tilde{J}(\theta) = \frac{1}{N}\sum_{i=1}^{N}\sum_{t=1}^{T} \log \pi_\theta(a_{i,t} \mid s_{i,t}) \cdot \text{adv}_{i,t}.$$
Esta forma facilita calcular el gradiente mediante diferenciación automática, y permite que cada dato participe simultáneamente en la optimización de varios pasos, sin necesidad de retropropagar por separado $N \times T$ veces.

\subsection{Resumen de la sección}

Esta sección analiza el gradiente de política desde una perspectiva intuitiva: el énfasis está en la imitación ponderada por reward, el reward-to-go que aporta causalidad, y el papel de baseline/advantage en la reducción de la variance. La función objetivo sustituta permite implementar el cálculo del gradiente de manera estándar en el marco.

